\documentclass[aps,prd,reprint,superscriptaddress,nofootinbib,floatfix]{revtex4-2}
\usepackage[T1]{fontenc}
\usepackage{mathpazo}
\usepackage{amsmath,amssymb}
\usepackage{graphicx}
\usepackage{xcolor}
\usepackage{microtype}
\usepackage[colorlinks=true,linkcolor=red!70!black,citecolor=blue!70!black,urlcolor=cyan!50!black]{hyperref}
\usepackage{orcidlink}
\usepackage{xurl}
\usepackage{adjustbox}
\renewcommand{\arraystretch}{1.12}
\newcommand{\DC}{\mathrm{DC}}
\graphicspath{{figures/}}

\begin{document}
\raggedbottom

\title{Dual Credit as a Readiness Lever: School-Level Evidence from Indiana\\and a Map of Where College Mathematics Pathways Are Missing}

\author{@@AUTHOR@@\,\orcidlink{@@ORCID@@}}
\email{@@EMAIL@@}
\affiliation{Independent Researcher, @@AFFIL@@}

\date{October 2026}

\begin{abstract}
States have expanded dual credit on the expectation that college coursework taken in high school prepares students for college, yet most evidence compares students who chose dual credit with students who did not, so the effect of a school's investment is hard to separate from who enrolls. We link Indiana Commission for Higher Education first-year college outcomes to federal Civil Rights Data Collection course offerings and NCES Common Core of Data records for @@NSCH@@ public high schools and @@NSC@@ school-cohorts (graduating classes 2021--2024; @@NINPUB@@ graduates who enrolled at an Indiana public college). A 10-percentage-point higher share of graduates earning dual credit is associated with a @@G1@@ higher school-average first-year GPA before adjustment. After accounting for school poverty, Advanced Placement success, size, locale, the timing of Algebra~I, and calculus and advanced-mathematics offerings, the association falls by @@ATTG@@ to @@G3@@ (@@G3CI@@); school fixed-effects estimates are similar (@@G5@@). First-year credits follow the same pattern. Because the state's first-year GPA and credit measures include dual-credit coursework recorded before enrollment, these estimates are upper bounds on any readiness effect; remediation, which the overlap does not touch, shows no robust association. Within schools, dual-credit earners out-perform classmates by @@WGAP@@ GPA points, a gap that widens where dual credit is common, consistent with strong selection. Access to the mathematics pathway is uneven: @@ANOCALC@@ of traditional public high schools offered no calculus in either of the two most recent federal collections, rising from @@Q1NC@@\% of the lowest-poverty quarter of schools to @@Q4NC@@\% of the highest, and in @@NCNOCALC@@ counties no public high school offered it. Any school-level readiness return to dual credit is modest at most; the larger policy problem is where the advanced mathematics pathway does not exist. All code, data and an interactive atlas are public.\\[4pt]
\textbf{Keywords:} dual credit, dual enrollment, college readiness, mathematics pathways, Algebra~I, Civil Rights Data Collection, Indiana
\end{abstract}

\maketitle

%=====================================================================
\section{Introduction}\label{sec:intro}

Dual credit---college coursework that high school students complete for both high school and college credit---has moved from an enrichment option for a few strong students to a core part of state college-readiness policy. Indiana has pushed further than most states. In the public high schools studied here, @@DCWMEAN@@ of the 2021--2024 graduates earned dual credit from an Indiana public college, and the median school's share was @@DCMED@@. The state's Indiana College Core, a 30-credit block of general-education courses that transfers to every Indiana public institution, was offered in 328 high schools in 2026--27, up from 84 five years earlier~\cite{cherelease}. In September 2026 the Commission for Higher Education (CHE) issued a Request for Information on reducing barriers to early-college access, noting that about one in five Hoosier students attends a high school without the College Core and attributing the gap mainly to low enrollment and a shortage of teachers credentialed to teach college courses~\cite{icc}.

These investments rest on two premises. The first is that dual credit prepares students for college: that a school where more students take college courses in high school sends graduates who perform better once they arrive. The second is that dual credit is available to students who would benefit from it. Both premises can be examined with public data, and both matter for where the next dollar goes.

Two features of the existing evidence limit what it can tell policymakers. First, most studies compare students who take dual credit with students who do not~\cite{an2013,karp2007,speroni2011,giani2014}. Even with rich controls, those comparisons mix the effect of the course with the characteristics of students who choose it and are admitted to it. The decision a state or district actually faces is different: whether to invest in dual-credit capacity at a school, not whether a particular student should enroll. Second, few studies ask whether schools that serve students with the least access to advanced coursework can offer dual credit at all, particularly in mathematics, where college-level work depends on a sequence that begins with Algebra~I.

This paper addresses both gaps with public data. It asks three questions:
\begin{enumerate}
\item Do Indiana high schools where more students earn dual credit send graduates who do better in their first year of college, after accounting for school poverty, Advanced Placement success, when students reach Algebra~I, and which mathematics courses the school offers?
\item Does any association differ for high-poverty, rural, or late-Algebra schools?
\item Where is the mathematics pathway into dual credit missing?
\end{enumerate}

We make four contributions. First, we construct a linked school-cohort panel joining CHE graduate outcomes to federal course-offering data for every Indiana public high school, and release it with full provenance. Second, we estimate the school-level association under a sequence of increasingly demanding specifications, including school fixed effects, and show how much of the raw advantage reflects which schools offer dual credit. Third, we formalize two reasons the raw pattern overstates any readiness return---selection of students within schools and the inclusion of dual-credit coursework in the state's first-year outcome measures---and bound the estimates accordingly. Fourth, we map, school by school and county by county, where the calculus pathway and dual enrollment are absent.

@@M118A@@

The remainder of the paper is organized as follows. Section~\ref{sec:background} reviews dual credit in Indiana and the related literature. Section~\ref{sec:data} describes the data and the construction of the panel, and Sec.~\ref{sec:measures} defines the measures. Section~\ref{sec:methods} sets out the empirical strategy. Section~\ref{sec:results} reports the readiness results and Sec.~\ref{sec:access} the access analysis. Section~\ref{sec:discussion} discusses implications, Sec.~\ref{sec:limitations} lists limitations, and Sec.~\ref{sec:conclusion} concludes.

%=====================================================================
\section{Background}\label{sec:background}

\subsection{Dual credit in Indiana}

Indiana delivers dual credit mainly through two routes. In the first, credentialed high school teachers teach college courses in their own high schools under a college's supervision; Indiana University's Advance College Project (ACP) and Ivy Tech Community College's high school programs are the largest examples. In the second, students take courses on a college campus or online. CHE reports, for every high school, how many graduates earned dual credit from an Indiana public college and how those graduates fared in their first college year~\cite{chedata}. CHE's own descriptive analysis of earlier cohorts found that graduates with pre-college credit had higher first-year GPAs, earned more credits, and needed less remediation than other students~\cite{che2016}. That comparison, like most in the field, is between students rather than schools.

\subsection{Evidence on dual enrollment and college outcomes}

Quasi-experimental studies generally find positive associations between dual enrollment and college enrollment, persistence, and degree attainment, though estimates shrink as controls for prior achievement tighten~\cite{an2013,karp2007,speroni2011}. Studies in Texas~\cite{giani2014,miller2017} and New York City~\cite{allen2012} report positive associations with college outcomes that vary with the type and amount of coursework, and a national analysis links dual enrollment to high school graduation, college enrollment, institutional choice, and persistence~\cite{lee2022}. The What Works Clearinghouse rated dual-enrollment programs as having positive effects on several college outcomes~\cite{wwc2017}. Work on equity documents that benefits and participation are unequally distributed: gaps by income and race persist even as programs expand~\cite{taylor2015,fink2017}.

\subsection{Selection and the school-level view}

Student-level comparisons face a basic problem. Students who take dual credit are recommended by teachers, meet grade or test thresholds, and choose to enroll; each of these predicts college success on its own. A school-level design changes the question. It relates a school's dual-credit uptake to the average outcomes of all its college-going graduates, so that selection of individual students into dual credit within a school does not by itself produce an association. Selection of schools into high uptake remains, which we address with controls and with school fixed effects that use only changes within a school over time.

\subsection{Algebra~I timing and advanced mathematics}

Students who take more advanced mathematics in high school do better in college~\cite{long2012}, but accelerating students into Algebra~I before they are ready can backfire: district policies that moved more students into eighth-grade algebra lowered later mathematics achievement for students at the margin~\cite{clotfelter2015,domina2015}. For dual-credit mathematics, the timing of Algebra~I matters because it determines whether a student can complete Algebra~II---the usual prerequisite for a college mathematics course---by eleventh or twelfth grade. We therefore control for Algebra~I timing and for whether a school offers calculus and other advanced mathematics.

\subsection{Remediation}

Students placed into remedial (developmental) mathematics or English in college often fail to complete the sequence and are less likely to earn a degree~\cite{bailey2010}. Remediation is therefore a natural readiness outcome. In Indiana, remediation is identified when an institution reports a student enrolled in remedial credits during the year after high school graduation~\cite{cheguide}, and it has become uncommon in recent cohorts (Fig.~\ref{fig:overview}).

%=====================================================================
\section{Data}\label{sec:data}

\subsection{Graduate outcomes}

CHE publishes an annual College-Going (formerly College Readiness) data set for each high school graduating cohort~\cite{chedata}. For each school it reports the number of graduates; the numbers enrolling in any college and in an Indiana public college; and, for Indiana public college enrollees, the number needing remediation (overall and by subject), average first-year GPA, and average first-year credits earned. Each measure is reported overall and separately for graduates who did and did not earn dual credit from an Indiana public college, by free/reduced-lunch status, and by Advanced Placement results. Small cells are suppressed. We use the 2021, 2022, 2023, and 2024 cohorts for the main analysis and the 2019 cohort as a pre-pandemic sensitivity check; the 2020 cohort, whose first college year began in fall 2020, is omitted.

According to CHE's user guide, first-year GPA is the student's latest GPA for the school year after high school graduation, and first-year credits are the cumulative credit hours earned during that year~\cite{cheguide}. Dual-credit grades are recorded by the sponsoring college---ACP instructors, for example, upload final grades to IU---and are on a student's record when the student enrolls. The GPA and credit measures therefore include dual-credit coursework taken in high school; Sec.~\ref{sec:overlap} examines the consequences.

\subsection{Course offerings and Algebra~I timing}

The U.S.\ Department of Education's Civil Rights Data Collection (CRDC) reports, for every public school, whether it offers dual enrollment and how many students participate; the number of sections and students in calculus and in advanced mathematics; and Algebra~I enrollment for grades 7--8, 9--10, and 11--12~\cite{crdc}. We use the 2021--22 collection for the 2021 and 2022 cohorts and the 2023--24 collection for the 2023 and 2024 cohorts; when a measure is missing in the matched year we use the other year. In Indiana, the 2021--22 collection includes @@CRDC22G12@@ schools offering grade 12 and the 2023--24 collection @@CRDC24G12@@.

\subsection{Directory, membership, and geography}

The NCES Common Core of Data (CCD) 2023--24 directory links CHE's state school numbers to NCES identifiers, and CCD membership files supply grade-8 and grade-9 enrollment as denominators for Algebra~I measures~\cite{ccd}. NCES EDGE geocodes supply coordinates and locale codes~\cite{edge}, and Census cartographic boundary files supply county outlines~\cite{census}.

\subsection{Linkage and sample}

CHE's 2024 file carries Indiana Department of Education (IDOE) school numbers; earlier files do not. We link earlier cohorts to IDOE numbers by exact school name through the 2024 file, which succeeds for @@LK19L@@\% of 2019 rows and nearly all later rows; unlinked rows are mainly nonpublic, closed, or renamed schools, each listed in the quality-assurance report released with the data. IDOE numbers are then matched to CCD, and through NCES identifiers to CRDC. Nonpublic schools (absent from CCD and CRDC) and Department of Correction facilities are dropped. Table~\ref{tab:sample} traces the sample. The analytic sample contains public schools with at least @@MININPUB@@ Indiana public enrollees in a cohort and non-missing outcomes, exposure, and controls: @@NSC@@ school-cohorts from @@NSCH@@ schools, covering @@NGRAD@@ graduates and @@NINPUB@@ Indiana public college enrollees.

@@T_SAMPLE@@

\subsection{Quality assurance}

The build script writes a quality-assurance report that records row counts at each stage, match rates, range checks (all shares in $[0,1]$, GPA in $[0,4]$, credits in $[0,60]$, dual-credit and non-dual-credit graduates summing to the total), and named spot checks of individual schools against the source files. All checks pass. Every number in this paper is read from a statistics file written by the analysis pipeline~\cite{repo}.

%=====================================================================
\section{Measures}\label{sec:measures}

Index schools by $j$ and graduating cohorts by $t$. Table~\ref{tab:vars} in Appendix~\ref{app:vars} defines every variable.

\paragraph*{Exposure.} The dual-credit share is
\begin{equation}
s_{jt}=\frac{N^{\DC}_{jt}}{N_{jt}},\label{eq:share}
\end{equation}
where $N_{jt}$ is the number of graduates and $N^{\DC}_{jt}$ the number who earned dual credit from an Indiana public college. When CHE suppresses one half of the two-way split, it is recovered as the total minus the other half. Coefficients are reported per 10 percentage points, i.e., for the regressor $10\,s_{jt}$.

\paragraph*{Outcomes.} For the $w_{jt}$ graduates of school $j$ in cohort $t$ who enrolled at an Indiana public college, CHE reports average first-year GPA and credits,
\begin{equation}
\bar G_{jt}=\frac{1}{w_{jt}}\sum_{i\in jt}G_i,\qquad \bar C_{jt}=\frac{1}{w_{jt}}\sum_{i\in jt}C_i,\label{eq:outcomes}
\end{equation}
and the share needing remediation, $\rho_{jt}=R_{jt}/w_{jt}$, where $R_{jt}$ counts enrollees placed in remedial coursework. Mathematics remediation (math only, or both math and English) accounts for @@MATHREMSHARE@@ of all remediation in the sample.

\paragraph*{Algebra~I timing.} CRDC reports seventh- and eighth-grade Algebra~I at middle schools, not at the high school a student later attends. The main measure is therefore defined for school corporation $c(j)$:
\begin{equation}
a^{\rm corp}_{c}=\frac{\sum_{k\in c}A^{(7\text{--}8)}_{k}}{\sum_{k\in c}M^{(8)}_{k}},\label{eq:algcorp}
\end{equation}
where $A^{(7\text{--}8)}_k$ is Algebra~I enrollment in grades 7--8 at school $k$ and $M^{(8)}_k$ is eighth-grade membership; it is computed for corporations with at least 20 eighth graders (@@CORPS22@@ corporations in 2021--22, median @@CORPMED22@@; @@CORPS24@@ in 2023--24, median @@CORPMED24@@) and capped at 1. A high-school alternative,
\begin{equation}
a^{\rm hs}_{j}=\frac{A^{(9\text{--}10)}_{j}+A^{(11\text{--}12)}_{j}}{M^{(9)}_{j}},\label{eq:alghs}
\end{equation}
captures late Algebra~I taking and is capped at 1.2 because repeaters can push it above one.

\paragraph*{Mathematics offerings and controls.} Indicators $k_j$ and $m_j$ equal one if the school ran at least one calculus or advanced-mathematics section. Controls include the free/reduced-lunch share of graduates $f_{jt}$, the share passing an AP exam $p_{jt}$, log graduates, and NCES locale $\ell_j$ (city, suburb, town, rural).

@@T_DESC@@

%=====================================================================
\section{Empirical strategy}\label{sec:methods}

\subsection{Pooled weighted regression}

The main specification relates school-cohort outcomes to dual-credit share:
\begin{equation}
Y_{jt}=\alpha+\beta\,(10\,s_{jt})+\mathbf{x}_{jt}^{\top}\boldsymbol{\gamma}+\delta_t+\varepsilon_{jt},\label{eq:main}
\end{equation}
where $Y_{jt}\in\{\bar G_{jt},\bar C_{jt},\rho_{jt}\}$, $\delta_t$ are cohort indicators, and $\mathbf{x}_{jt}$ collects controls added in stages: M1 has none; M2 adds $f_{jt}$, $p_{jt}$, log graduates, and locale; M3, the preferred specification, adds $a^{\rm corp}_{c(j)}$, $k_j$, and $m_j$; M4 adds the in-state public college-going rate $w_{jt}/N_{jt}$. Because outcomes are averages over $w_{jt}$ enrollees, the model is estimated by weighted least squares,
\begin{equation}
\hat{\boldsymbol\theta}=\arg\min_{\boldsymbol\theta}\sum_{j,t}w_{jt}\bigl(Y_{jt}-\mathbf{z}_{jt}^{\top}\boldsymbol\theta\bigr)^2,\label{eq:wls}
\end{equation}
with $\mathbf{z}_{jt}$ the full regressor vector.

\subsection{Inference}

School-cohorts from the same school are not independent. Standard errors are cluster-robust by school,
\begin{equation}
\widehat{\mathbf V}=\frac{G}{G-1}\frac{N-1}{N-K}\,\mathbf B^{-1}\Bigl(\sum_{g=1}^{G}\mathbf Z_g^{\top}\mathbf W_g\hat{\mathbf e}_g\hat{\mathbf e}_g^{\top}\mathbf W_g\mathbf Z_g\Bigr)\mathbf B^{-1},\label{eq:cluster}
\end{equation}
with $\mathbf B=\mathbf Z^{\top}\mathbf W\mathbf Z$, $G=@@NSCH@@$ clusters, $N=@@NSC@@$ observations, and $K$ regressors.

\subsection{School fixed effects}

Stable differences between schools---in preparation, culture, or the colleges graduates attend---may drive both uptake and outcomes. Model M5 replaces time-invariant controls with school effects $\mu_j$,
\begin{equation}
Y_{jt}=\mu_j+\beta\,(10\,s_{jt})+\delta_t+\varepsilon_{jt},\label{eq:fe}
\end{equation}
so that $\beta$ is identified from year-to-year changes in a school's uptake. Those changes are small, so M5 is imprecise.

\subsection{Bounding by college-going}

Outcomes are observed only for graduates who enroll at an Indiana public college. If dual credit changes who enrolls---for example, drawing marginal students into college---average outcomes can move for reasons unrelated to readiness. M4 conditions on the in-state public college-going rate. Because that rate may itself respond to dual credit, M4 is read as a bound rather than the preferred estimate.

\subsection{Heterogeneity}

To ask for which schools the association is larger, we interact dual-credit share with an indicator $D_j$ (high poverty, rural, below-median early Algebra~I, or calculus offered):
\begin{equation}
Y_{jt}=\dots+\beta\,(10\,s_{jt})+\theta\,(10\,s_{jt})D_j+\kappa D_j+\varepsilon_{jt}.\label{eq:het}
\end{equation}
The slope in the group is $\beta+\theta$, with variance $\widehat{\rm Var}(\hat\beta)+\widehat{\rm Var}(\hat\theta)+2\,\widehat{\rm Cov}(\hat\beta,\hat\theta)$.

\subsection{Within-school gaps and selection}

Where CHE reports both groups, the within-school gap is
\begin{equation}
\Delta_{jt}=\bar G^{\DC}_{jt}-\bar G^{\rm no}_{jt},\label{eq:gap}
\end{equation}
and the school average is the enrollee-weighted mix $\bar G_{jt}=\tilde s_{jt}\bar G^{\DC}_{jt}+(1-\tilde s_{jt})\bar G^{\rm no}_{jt}$, where $\tilde s_{jt}$ is the dual-credit share among enrollees. If dual credit raised every participant's outcomes by the same amount and participants were a random draw, $\Delta_{jt}$ would not depend on $s_{jt}$. If instead schools enroll their strongest students first, expanding a program moves progressively less-prepared students out of the non-participant group and leaves the weakest behind, so that $\bar G^{\rm no}_{jt}$ falls and $\Delta_{jt}$ rises with $s_{jt}$. We estimate
\begin{equation}
\Delta_{jt}=\alpha+\lambda\,(10\,s_{jt})+\phi\,(10f_{jt})+\delta_t+e_{jt},\label{eq:gapreg}
\end{equation}
weighted by the number of enrollees in both groups. A positive $\lambda$ is evidence of selection.

\subsection{Overlap between exposure and outcome}\label{sec:overlap}

Let student $i$ carry $h_i$ dual-credit hours with grade average $g^{\DC}_i$ onto the college record, and earn $H_i$ hours with average $g^{\rm col}_i$ in college courses during the first year. If the reported measures include the dual-credit record,
\begin{align}
C_i&=H_i+\pi h_i,\label{eq:credmech}\\
G_i&=\frac{H_i\,g^{\rm col}_i+\pi h_i\,g^{\DC}_i}{H_i+\pi h_i},\label{eq:gpamech}
\end{align}
where $\pi\in[0,1]$ is the fraction of dual-credit hours that appear in the institution's first-year record (one if all do; less if, for example, only credit from the same institution is counted). Averaging Eq.~(\ref{eq:credmech}) within a school gives $\bar C_{jt}=\bar H_{jt}+\pi\,\tilde s_{jt}\,\bar h^{\DC}_{jt}$, where $\bar h^{\DC}$ is mean dual-credit hours among earners. A 10-point rise in $\tilde s$ therefore raises reported credits mechanically by about $0.1\,\pi\,\bar h^{\DC}$ even if college performance is unchanged. CHE reported that 2014 graduates who earned dual credit averaged about 10 dual-credit hours~\cite{che2016}; with $\pi=1$ the mechanical term alone would be on the order of @@MECHC@@ credit per 10 points. The unadjusted credits slope, @@C1@@ per 10 points (Table~\ref{tab:main}), is smaller than that, which suggests $\pi<1$ in practice, but the adjusted estimates cannot be separated from the mechanical term with published data. For GPA, Eq.~(\ref{eq:gpamech}) blends dual-credit grades into the first-year average, so schools with more earners will report different first-year GPAs whenever $g^{\DC}\neq g^{\rm col}$. We therefore read all GPA and credit estimates as upper bounds on any effect of dual credit on first-year college performance. The remediation outcome $\rho_{jt}$ is based on placement after enrollment and is not affected by this overlap.

%=====================================================================
\section{Results: dual credit and first-year outcomes}\label{sec:results}

\subsection{Descriptive patterns}

Figure~\ref{fig:overview} summarizes the sample. Dual-credit uptake is high almost everywhere and has been stable: the graduate-weighted share rose from @@TR19DC@@ for the 2019 cohort to @@TR24DC@@ for 2024. Average first-year GPA of in-state public enrollees moved from @@TR21GPA@@ (2021) to @@TR24GPA@@ (2024), and the share needing remediation fell from @@TR19REM@@ (2019) to @@TR21REM@@ (2021) and @@TR24REM@@ (2024). The in-state public college-going rate fell from @@TR19INP@@ to @@TR24INP@@ over the same period.

\begin{figure*}[t]
\includegraphics[width=\textwidth]{fig1_overview.pdf}
\caption{Overview of the analytic sample. Top: key figures. Bottom: dual-credit share of graduates, first-year GPA, and the share of enrollees needing remediation, by graduating cohort (2019 shown for reference; 2020 omitted). Cohort aggregates use all linked public school-cohorts with at least @@MININPUB@@ in-state public enrollees.}
\label{fig:overview}
\end{figure*}

Across school-cohorts the dual-credit share is concentrated between @@DCQ25@@ and @@DCQ75@@ (interquartile range), with a long lower tail (Fig.~\ref{fig:dist}); @@NDCZERO@@ school-cohorts report no dual-credit earners and @@NDC90@@ report 90\% or more. Higher-uptake schools are less poor (correlation with free/reduced-lunch share @@CORRDCFRL@@), smaller, and more often rural (Table~\ref{tab:terc}). The raw correlation between dual-credit share and first-year GPA is @@CORRDCGPA@@.

\begin{figure}[t]
\includegraphics[width=\columnwidth]{fig2_dc_distribution.pdf}
\caption{Distribution of school dual-credit share across the @@NSC@@ school-cohorts of the analytic sample.}
\label{fig:dist}
\end{figure}

@@T_TERC@@

\subsection{Main estimates}

Table~\ref{tab:main} and Fig.~\ref{fig:coef} report $\hat\beta$ from Eq.~(\ref{eq:main}). Without controls, a 10-point higher dual-credit share is associated with @@G1@@ GPA points and @@C1@@ first-year credits. Adding poverty, AP success, size, and locale reduces the GPA estimate to @@G2@@ (@@G2CI@@) and raises $R^2$ from @@R2M1@@ to @@R2M3@@. Adding Algebra~I timing and mathematics offerings leaves it essentially unchanged at @@G3@@ (@@G3CI@@, @@G3P@@). Neither early Algebra~I (@@ALGB@@ per 0.1, @@ALGP@@) nor calculus (@@CALCB@@, @@CALCP@@) nor advanced mathematics (@@ADVMB@@, @@ADVMP@@) predicts school first-year GPA once poverty (@@FRLB@@ per 10 points) and AP success (@@APB@@ per 10 points) are held fixed (Table~\ref{tab:full}). Defining attenuation as $1-\hat\beta_{\rm M3}/\hat\beta_{\rm M1}$, school characteristics account for @@ATTG@@ of the raw GPA association and @@ATTC@@ of the raw credits association.

\begin{figure*}[t]
\includegraphics[width=\textwidth]{fig3_coefficients.pdf}
\caption{Estimated change in school-average first-year GPA (left) and credits (right) per 10-point higher dual-credit share, by specification. Whiskers are 95\% confidence intervals from Eq.~(\ref{eq:cluster}). M3 (navy) is the preferred specification.}
\label{fig:coef}
\end{figure*}

@@T_MAIN@@

The school fixed-effects estimate, @@G5@@ (@@G5CI@@, @@G5P@@), is close to M3, which suggests that stable unobserved school differences do not drive the adjusted estimate, although its interval is wide. Conditioning on in-state public college-going (M4) lowers the GPA estimate to @@G4@@ (@@G4P@@) and the credits estimate to @@C4@@ (@@C4P@@). Figure~\ref{fig:binscatter} shows the same attenuation visually: once both axes are residualized on the M3 controls, the steep raw gradient flattens.

\begin{figure}[t]
\includegraphics[width=\columnwidth]{fig4_binscatter.pdf}
\caption{School first-year GPA by dual-credit share, unadjusted and adjusted for M3 controls (20 equal-count bins; enrollee-weighted).}
\label{fig:binscatter}
\end{figure}

\subsection{Magnitudes and the overlap bound}

The 10th- and 90th-percentile schools differ by @@P1090@@ points of dual-credit share. Under M3 that difference implies at most about @@P1090G@@ GPA points and @@P1090C@@ credits, including any dual-credit coursework counted in the first-year record (Sec.~\ref{sec:overlap}). For comparison, the enrollee-weighted mean first-year GPA is @@GPAMEAN@@ and mean credits are @@CRMEAN@@. These are upper bounds; the true difference in first-year college performance attributable to dual credit could be smaller.

\subsection{Remediation}

The share needing remediation shows a small negative raw association (@@R1@@ per 10 points, @@R1P@@) that disappears with controls (@@R3@@, @@R3P@@) and in the fixed-effects model (@@R5@@, @@R5P@@). Because remediation is measured after enrollment, it is the one outcome free of the overlap. Its null result may also reflect how rare remediation has become: the enrollee-weighted mean is @@REMMEAN@@.

\subsection{For which schools}

Point estimates are larger in high-poverty schools (@@HFRLG@@ GPA per 10 points) than in lower-poverty schools (@@HFRLO@@), and likewise for credits (@@HFRLCG@@ versus @@HFRLCO@@), but neither difference is statistically distinguishable from zero (@@HFRLP@@ and @@HFRLCP@@). There is no evidence of different slopes for rural schools (@@HRURP@@), schools in corporations with less early Algebra~I (@@HALGP@@), or schools offering calculus (@@HCALCP@@) (Table~\ref{tab:het}, Fig.~\ref{fig:het}). The data are consistent with a modest, broadly shared association and cannot rule out moderately larger returns in high-poverty schools.

\begin{figure}[t]
\includegraphics[width=\columnwidth]{fig6_heterogeneity.pdf}
\caption{Group-specific GPA slopes from Eq.~(\ref{eq:het}). Dark bars: the named group; light bars: the complement. The $p$-value tests equality of the two slopes.}
\label{fig:het}
\end{figure}

@@T_HET@@

\subsection{Within schools: selection}

In the @@WSCH@@ schools (@@WSC@@ school-cohorts) where both groups are reported, graduates who earned dual credit averaged @@WGAP@@ GPA points and @@WCGAP@@ credits more in their first college year than classmates who did not; the gap was positive in @@WPOS@@ of school-cohorts. Estimating Eq.~(\ref{eq:gapreg}), the gap widens by @@WSLOPE@@ GPA points per 10 points of dual-credit share (@@WSLOPEP@@), and the non-earners' average falls by @@WNODC@@ per 10 points, conditional on school poverty (Table~\ref{tab:within}, Fig.~\ref{fig:within}). This is the pattern Sec.~\ref{sec:methods} predicts under selection: as a program grows, the students it does not reach are increasingly the least prepared. Part of the earners' advantage is also mechanical, since their first-year records include their dual-credit grades and hours. Student-level comparisons of earners and non-earners should therefore be read as upper bounds.

\begin{figure}[t]
\includegraphics[width=\columnwidth]{fig5_within_school.pdf}
\caption{First-year GPA of dual-credit earners and non-earners within the same schools, by sextile of school dual-credit share.}
\label{fig:within}
\end{figure}

@@T_WITHIN@@

\subsection{Sensitivity}

Across the checks in Table~\ref{tab:sens} and Fig.~\ref{fig:sens}, the adjusted GPA estimate ranges from @@SENSGLO@@ to @@SENSGHI@@ and the credits estimate from @@SENSCLO@@ to @@SENSCHI@@. The GPA estimate is larger without weights (@@SUNW@@), indicating that smaller schools contribute more to the association, and smaller and not statistically significant when restricted to school-cohorts with 30 or more in-state enrollees (@@SMIN30@@, @@SMIN30P@@). Replacing the corporation Algebra~I measure with the high-school measure (@@SHSALG@@), including both, or adding the 2019 cohort (@@S2019@@) changes little; the 2024 cohort alone gives @@S2024@@ (@@S2024P@@). The credits estimate is less stable: it falls to @@SNOVIRTC@@ (@@SNOVIRTCP@@) when virtual schools are excluded, while the GPA estimate is unchanged (@@SNOVIRTG@@).

\begin{figure*}[t]
\includegraphics[width=\textwidth]{fig7_sensitivity.pdf}
\caption{Sensitivity of the M3 estimate for first-year GPA (left) and credits (right). Whiskers are 95\% confidence intervals.}
\label{fig:sens}
\end{figure*}

@@T_SENS@@

%=====================================================================
\section{Where the mathematics pathway is missing}\label{sec:access}

\subsection{Definitions}

Neither CHE nor CRDC identifies the subject of dual-credit courses, so dual-credit mathematics cannot be observed directly. We approximate the pathway with two observable conditions: whether a school offers calculus, the end point of the sequence that feeds college-level mathematics, and whether it offers dual enrollment at all, together with its graduates' dual-credit uptake. Calculus and dual enrollment count as offered if reported in either CRDC 2021--22 or 2023--24, because small schools often run calculus in alternate years; calculus status changed between the two collections at @@CALCCHANGED@@ schools. For traditional public high school $j$ with calculus indicator $k_j$, dual-enrollment indicator $e_j$, and pooled 2021--2024 dual-credit share $\bar s_j$, the access tier is
\begin{equation}
\tau_j=\begin{cases}
\text{No dual enrollment} & e_j=0,\\
\text{Thin pathway} & e_j=1,\ k_j=0,\ \bar s_j<\tfrac13,\\
\text{Dual credit, no calculus} & e_j=1,\ k_j=0,\ \bar s_j\geq\tfrac13,\\
\text{Dual credit and calculus} & e_j=1,\ k_j=1.
\end{cases}\label{eq:tier}
\end{equation}
Virtual and alternative schools (@@AVALT@@) are excluded.

\subsection{Statewide}

Of @@ATRAD@@ traditional public high schools matched to CRDC 2023--24, @@ANOCALCN@@ (@@ANOCALC@@) offered no calculus in either year; they enroll @@ANOCALCSTN@@ students, or @@ANOCALCST@@ of the @@ATRADST@@ students in these schools (Table~\ref{tab:tiers}). @@TNONE@@ schools reported no dual enrollment, and a further @@TTHIN@@ had neither calculus nor substantial dual-credit uptake. Schools without calculus are smaller (median enrollment @@NCENR@@ versus @@CENR@@), poorer (median free/reduced-lunch share @@NCFRL@@ versus @@CFRL@@), and send fewer graduates with dual credit (median share @@NCDC@@ versus @@CDC@@). Most still offer some dual enrollment: @@NOCALCDUAL@@ of the @@ANOCALCN@@ do.

@@T_TIERS@@

\subsection{Poverty and locale}

Absence of calculus tracks poverty closely: @@Q1NC@@\% of the lowest-poverty quarter of schools lacked it, rising to @@Q2NC@@\%, @@Q3NC@@\%, and @@Q4NC@@\% across the higher quartiles (Fig.~\ref{fig:accesscols}). Median dual-credit uptake falls from @@Q1DC@@ in the lowest-poverty quartile to @@Q4DC@@ in the highest. By locale, the gap is largest in cities: @@CITYNC@@\% of city high schools offered no calculus and @@CITYND@@\% no dual enrollment, against @@RURNC@@\% and @@RURND@@\% of rural schools (Table~\ref{tab:access}, Fig.~\ref{fig:tiers}). In Marion County alone, @@MARIONNC@@ of @@MARION@@ traditional public high schools lacked calculus. By students, @@CITYTHIN@@\% of city high school students attend a school with no dual enrollment or a thin pathway, compared with @@RURTHIN@@\% in rural areas.

\begin{figure*}[t]
\includegraphics[width=\textwidth]{fig8_access_columns.pdf}
\caption{Share of traditional public high schools with no calculus (in either CRDC year) and with no dual enrollment, by school poverty quartile (left) and NCES locale (right).}
\label{fig:accesscols}
\end{figure*}

\begin{figure}[t]
\includegraphics[width=\columnwidth]{fig9_tiers.pdf}
\caption{Distribution of traditional public high schools across the access tiers of Eq.~(\ref{eq:tier}), by locale.}
\label{fig:tiers}
\end{figure}

@@T_ACCESS@@

\subsection{Geography}

Figure~\ref{fig:map} maps, for each county, the share of public high school students in schools with no calculus, with schools reporting no dual enrollment or a thin pathway marked. In @@NCNOCALC@@ counties (@@CNOCALCLIST@@) no public high school offered calculus in either year, and in @@NCMAJ@@ counties most students attend a school without it (Table~\ref{tab:county}). Urban absences concentrate in Indianapolis and the Gary--Hammond area. Each of the counties with no calculus anywhere has a single public high school, so one school's course list decides access for the whole county. An interactive version of the map with a searchable school register accompanies the replication package~\cite{repo}.

\begin{figure}[t]
\includegraphics[width=\columnwidth]{fig10_map.pdf}
\caption{Share of public high school students in schools without calculus, by county, with schools reporting no dual enrollment (circles) or a thin pathway (diamonds).}
\label{fig:map}
\end{figure}

@@T_COUNTY@@

%=====================================================================
\section{Discussion}\label{sec:discussion}

\subsection{What the readiness estimates mean}

The school-level association between dual credit and first-year college outcomes in Indiana is positive, survives adjustment for poverty, AP success, Algebra~I timing, and mathematics offerings, and appears in school fixed-effects models. It is also modest, most of the raw gap that motivates dual-credit advocacy reflects differences between schools, and part of what remains reflects dual-credit grades and hours counted in the first-year record itself. The one outcome free of that overlap, remediation, shows no robust association. Within schools, the dual-credit advantage widens as programs grow, a sign that students who are not reached differ sharply from those who are.

Three implications follow for policy. First, statewide dual-credit expansion should be evaluated against realistic expectations: a large increase in uptake corresponds to a few hundredths of a GPA point and under one credit at the school level, and some of that is the dual-credit work itself. Second, the binding constraint for college mathematics in many schools is the absence of a pathway. A fifth of traditional public high schools, and four in ten of the highest-poverty schools, offered no calculus in either recent year. Dual-credit mathematics courses such as finite mathematics or college algebra can be offered without calculus, but they still require students who reach Algebra~I early enough and teachers credentialed to teach college courses. Third, where uptake is already high, the remaining gains lie with the students not yet reached, whose first-year outcomes are weakest.

\subsection{A view from the classroom}

@@M118B@@

\subsection{Measurement the state could add}

Two changes to public reporting would sharpen this analysis considerably. CHE could report school-level first-year outcomes that exclude credit earned before enrollment, which would remove the overlap of Sec.~\ref{sec:overlap}; and it could report dual credit by subject, at minimum separating mathematics and English, so that the state can see where college-level mathematics is reached. A first-year gateway-mathematics completion measure by high school would be the readiness outcome closest to the question this paper asks.

%=====================================================================
\section{Limitations}\label{sec:limitations}

\begin{enumerate}
\item \emph{Design.} Estimates are school-level and observational; they do not identify the effect of dual credit on any individual student, and schools that expand dual credit may differ in ways the controls and fixed effects do not capture.
\item \emph{Outcome population.} Outcomes cover only graduates who enroll at an Indiana public college the following year, about a third of graduates at the median school; enrollees out of state or at private colleges are not observed.
\item \emph{Overlap.} First-year GPA and credits include dual-credit coursework recorded before enrollment, so GPA and credit estimates are upper bounds (Sec.~\ref{sec:overlap}).
\item \emph{Subject and volume.} CHE records whether a graduate earned any dual credit from an Indiana public college, not the subject or number of hours, and excludes credit from private or out-of-state colleges. The access analysis uses calculus and overall uptake as proxies for dual-credit mathematics.
\item \emph{Algebra~I timing.} The corporation measure does not follow each graduating cohort, and the high-school measure is affected by repeaters.
\item \emph{CRDC.} Course data are self-reported by schools and sometimes inconsistent between collections.
\item \emph{Small schools.} The association is concentrated in smaller schools; it is not statistically significant among schools with 30 or more in-state enrollees.
\item \emph{Pandemic.} The 2021--2023 cohorts completed high school during pandemic disruption; results are similar when the 2019 cohort is added.
\end{enumerate}

%=====================================================================
\section{Conclusion}\label{sec:conclusion}

Indiana's investment in dual credit coincides with modestly better first-year college records at the school level, after accounting for the schools' other advantages, though part of that difference is the dual-credit work itself carried onto the college transcript, and the readiness outcome untouched by that overlap shows no clear relationship. The clearer finding concerns access: the mathematics courses that lead into college-level mathematics are missing from many of the state's highest-poverty high schools, from four in ten city high schools, and from every public high school in @@NCNOCALC@@ counties. Policy aimed at readiness should measure dual credit by subject, report first-year outcomes that exclude credit earned before enrollment, and direct capacity---credentialed teachers, shared sections, earlier Algebra~I---to the schools where the pathway does not yet exist.

\section*{Data and code availability}
All data are public. Code, processed data, figures, the interactive atlas, and a full provenance log (source URLs, file sizes, and SHA-256 hashes) are available at \url{@@REPO@@}~\cite{repo}. Analyses used Python with pandas~\cite{mckinney2010} and statsmodels~\cite{seabold2010}; figures used matplotlib~\cite{hunter2007}.

\section*{Competing interests}
@@COMPETING@@

\section*{Funding}
@@FUNDING@@

\section*{Use of AI tools}
@@AIUSE@@

\section*{Author contributions}
@@CREDIT@@

%=====================================================================
\appendix

\section{Variable definitions}\label{app:vars}

Table~\ref{tab:vars} lists the symbols used in Secs.~\ref{sec:measures}--\ref{sec:methods}, their definitions, and their sources.

@@T_VARS@@

\section{Full coefficient estimates}\label{app:full}

Table~\ref{tab:full} reports all coefficients for the preferred specification and the college-going bound. Poverty and AP success are the strongest predictors of school first-year GPA; mathematics offerings and Algebra~I timing add little once they are included. Larger schools have slightly lower remediation rates.

@@T_FULL@@

\section{Reproduction}\label{app:repro}

The pipeline runs in numbered steps: \texttt{01\_fetch.py} downloads every source and logs its size and hash; \texttt{02\_build.py} links the files and writes the panel and quality-assurance report; \texttt{03\_analysis.py} estimates all models and writes the statistics file; \texttt{08\_preprint\_stats.py}, \texttt{09\_pbi\_figures.py}, and \texttt{10\_preprint.py} produce this paper. Raw downloads total about 1.4~GB; every later step runs offline, and a rerun from the logged sources reproduces the processed data and statistics file byte for byte.

%%BIBLIOGRAPHY%%

\end{document}
